\section{Discussion}
\label{sec:discussion}

\subsection{Key Findings}

Our experiments reveal that model rankings are highly stable between ImageNet and ImageNet-V2 ($\tau = 0.9647$), despite substantial accuracy degradation ($\sim$12\% mean drop). This finding has both theoretical and practical implications.

\subsubsection{Uniform Degradation Mechanism}

The high ranking correlation is explained by \emph{uniform} accuracy degradation. All models---regardless of architecture, parameter count, or publication year---experience similar percentage drops when evaluated on ImageNet-V2. This uniformity preserves ordinal relationships: if Model A beats Model B on ImageNet by 2\%, it typically still beats Model B on V2 by approximately 2\%.

\subsection{Limitations}

\textbf{Only One Benchmark Pair Tested.} Our analysis is limited to ImageNet $\rightarrow$ ImageNet-V2. Other distribution shifts (ObjectNet, ImageNet-Sketch, ImageNet-R) may show different patterns. Importantly, ImageNet-V2 was constructed to replicate the original data collection methodology, making it a ``near'' shift---the observed ranking stability may partially reflect this design choice rather than inherent model robustness.

\textbf{Sample May Not Represent Full Population.} We analyze 96 models with reported results on both benchmarks. Selection bias may exist---models expected to generalize well may be overrepresented.

\textbf{Original Hypothesis Refuted.} We hypothesized $\tau < 0.90$ and found $\tau = 0.96$. While this is a null result, it is scientifically valuable: it establishes that ranking instability is \emph{not} a significant concern for ImageNet-V2.

\subsection{Broader Impact}

This work has generally positive implications. Validating benchmark-based model selection reduces the risk of practitioners choosing suboptimal models due to misleading benchmarks. The finding does not enable harmful applications.
